% Industry CV
\documentclass[11pt,a4paper]{article}

\usepackage[utf8]{inputenc}
\usepackage[T1]{fontenc}
\usepackage{lmodern}
\usepackage[margin=1.6cm]{geometry}
\usepackage{xcolor}
\usepackage{hyperref}
\usepackage{enumitem}
\usepackage{tabularx}
\usepackage{fancyhdr}
\usepackage{lastpage}
\usepackage{microtype}
\usepackage{needspace}

\definecolor{accent}{HTML}{17365D}
\definecolor{muted}{HTML}{4A4A4A}

\hypersetup{
  colorlinks=true,
  urlcolor=accent,
  linkcolor=accent,
  pdfauthor={Antonio Michele Miti},
  pdftitle={Antonio Michele Miti - Industry CV}
}

\pagestyle{fancy}
\fancyhf{}
\renewcommand{\headrulewidth}{0pt}
\cfoot{\footnotesize\textcolor{muted}{Antonio Michele Miti \enspace | \enspace Page \thepage\ of \pageref{LastPage}}}
\setlength{\footskip}{18pt}

\setlength{\parindent}{0pt}
\setlength{\parskip}{0pt}
\setlist[itemize]{leftmargin=1.3em,itemsep=2.5pt,topsep=3pt,parsep=0pt,partopsep=0pt}

\newcommand{\cvsection}[1]{%
  \vspace{0.85em}%
  {\large\bfseries\color{accent} #1}\par
  \vspace{0.18em}\hrule\vspace{0.45em}%
}

\newcommand{\role}[4]{%
  \Needspace{6\baselineskip}%
  \textbf{#1}\hfill\textbf{#2}\par
  \textit{#3}\par
  \begin{itemize}#4\end{itemize}
  \vspace{0.3em}%
}

\newcommand{\project}[3]{%
  \textbf{#1}\hfill\textbf{#2}\par
  #3\par\vspace{0.7em}%
}

\begin{document}

{\Huge\bfseries Antonio Michele Miti}\par
\vspace{0.15em}
{\large\color{accent} Mathematician | Scientific Computing and Software Development}\par
\vspace{0.4em}
\small
\href{mailto:antoniomichele.miti@gmail.com}{antoniomichele.miti@gmail.com}
\enspace | \enspace +39 333 9825 418
\enspace | \enspace \href{https://www.antoniomiti.it/}{antoniomiti.it}
\enspace | \enspace \href{https://github.com/MasterToninus}{github.com/MasterToninus}\par
Italian citizen \enspace | \enspace Based in Lyon, France \enspace | \enspace Open to relocation or remote work

\cvsection{Profile}

Mathematician with an international research background and a longstanding interest in scientific programming. My work has trained me to organise difficult problems, make assumptions explicit, and develop arguments that can be checked in detail. Alongside research, I have developed a practical background in programming and Linux systems, which I am now looking to deepen in a professional software environment.

\cvsection{Technical Background}

\begin{tabularx}{\textwidth}{@{}>{\bfseries}p{4.6cm}>{\raggedright\arraybackslash}X@{}}
Programming languages & C++, CUDA, Haskell, Java, Mathematica, Python \\
Scientific computing & Algorithm implementation and optimisation, numerical experiments, quantum-circuit design \\
Systems and tools & Linux administration, Git, computer-cluster maintenance, operating-system installation and configuration \\
Research practice & Mathematical modelling, proof-based validation, technical writing, peer review, teaching, international collaboration
\end{tabularx}

\cvsection{Selected Technical Work}

\project{Quantum Algorithms for Modular Arithmetic}{DAC 2025}{
Co-developed a family of quantum circuits for arithmetic operations and proved the corresponding constructions. The proposed measurement-based uncomputation method reduced Toffoli count and depth by 10--25\%, depending on the circuit architecture. \href{https://arxiv.org/abs/2407.20167}{Paper and technical details}.}

\project{Functional Programming and Domain-Specific Languages}{KU Leuven, 2019}{
Attended a two-day industry course combining lectures with practical Haskell exercises. The programme covered higher-order functions, recursive programming, algebraic data types, pattern matching, and the construction of small embedded DSLs.}

\project{Linux Systems Administration}{Universit\`a Cattolica and LCM, 2014--2018}{
Maintained departmental Linux systems and a teaching computer cluster. The work included operating-system installation and configuration, user support, and assistance during computational-physics courses.}

\project{GPU Algorithm Implementation and Optimisation}{C++ / CUDA}{
Implemented a hybrid C++/CUDA version of a condensation algorithm for determinant computation. Analysed memory-access patterns, compared GPU kernels with a sequential CPU implementation across matrix sizes, and investigated the trade-off between performance and numerical stability.}

\cvsection{Professional Experience}

\role{Postdoctoral Fellow, CIVIS3i MSCA COFUND}{02/2024--03/2026}{Sapienza Universit\`a di Roma, Italy}{
  \item Led the MultiSymORI research project on multisymplectic geometry, observables, reduction, and their connections with mathematical physics.
  \item Managed the scientific work, international collaborations, research visits, reporting, publications, and conference presentations associated with the fellowship.
}


\role{Postdoctoral Researcher}{10/2022--01/2024}{Universit\`a Cattolica del Sacro Cuore, Brescia, Italy}{
  \item Worked on symplectic and multisymplectic geometry, collaborated on international research projects, and taught graduate-level mathematics.
}

\role{Research Fellow}{10/2021--09/2022}{Max Planck Institute for Mathematics, Bonn, Germany}{
  \item Studied higher algebraic structures in geometry and mathematical physics within an international research institute.
}

\role{Doctoral Researcher}{01/2017--04/2021}{KU Leuven, Belgium, and Universit\`a Cattolica del Sacro Cuore, Italy}{
  \item Completed a joint PhD on homotopy comomentum maps in multisymplectic geometry, spending extended periods at both institutions.
  \item Combined research with graduate teaching and departmental IT support.
}

\role{Java Developer Intern}{03/2016--07/2016}{CST s.r.l., Milan, Italy}{
  \item Trained in Java, object-oriented programming, and professional software-development tools and workflows.
}

\role{Computer Cluster and Linux Support Assistant}{05/2014--02/2015}{LCM, Universit\`a degli Studi di Milano, Italy}{
  \item Maintained the laboratory computer cluster and supported Linux installation, configuration, and computational-physics teaching.
}

\cvsection{Education}

\textbf{PhD in Mathematics}, Universit\`a Cattolica del Sacro Cuore and KU Leuven \hfill 2021\par
Joint international doctoral programme; thesis in differential geometry and mathematical physics.\par\vspace{0.25em}

\textbf{MSc in Theoretical Physics}, Universit\`a degli Studi di Milano \hfill 2015\par
Mathematical methods in physics, geometric mechanics, and computational methods.\par\vspace{0.25em}

\textbf{BSc in Physics}, Universit\`a degli Studi di Milano-Bicocca \hfill 2010

\cvsection{Selected Technical Training}

\begin{tabularx}{\textwidth}{@{}Xr@{}}
Functional Programming and Domain-Specific Languages, KU Leuven & 2019 \\
Scientific Computing with Python, Universit\`a Cattolica & 2018 \\
Python for Computational Science, CINECA-SCAI & 2017 \\
Introduction to Scientific and Technical Computing in C++, CINECA-SCAI & 2016 \\
Scientific Programming with GPGPU and CUDA, CINECA-SCAI & 2014
\end{tabularx}

\cvsection{Additional Information}

\textbf{Languages:} Italian (native), English (professional proficiency; TOEFL iBT 102/120), French (basic).\par
\textbf{Research output:} 8 scientific publications, including work on quantum algorithms; 20+ conference talks; experience reviewing technical research and mentoring students.

\end{document}
